\chapter{GENERAL CONCEPT AND SYSTEM COMPONENTS}

\section{Introduction}

Building on the general framework and objectives outlined in the previous chapter,
this chapter focuses on the technical environment that supports the development of the
IO‑Box project. It presents the hardware and software components employed, the
development tools and programming environment, as well as the communication
protocols that enable interaction with sensors, actuators, and external systems. The
chapter also introduces the hardware architecture, detailing how the different elements
are organized and integrated to form a coherent and functional platform.

By examining these aspects, we establish the practical foundation necessary for the
design and implementation of the IO‑Box. This technical overview highlights the
resources and methodologies that will be used throughout the project, ensuring that the
system is both reliable and adaptable to industrial requirements.

\section{Synoptic Diagram of the Proposed Solution}

The proposed solution can be represented by a synoptic diagram that highlights the
main functional phases of the system. The process begins with the acquisition of data
from the integrated sensors. These raw signals are then subjected to filtering and
calibration in order to remove noise and ensure measurement accuracy. The processed
data is passed to the regulation stage, where decision‑making and control logic are
applied. Finally, the regulated information leads to two possible actions: either the
transmission of data to external systems (IoT/cloud or monitoring dashboard), or the
activation of actuators to directly influence the physical environment.

\begin{figure}[H]
    \centering
    \includegraphics[width=0.75\textwidth]{C:/Users/selim/Downloads/imagesreport/synoptique.png}
    \caption{Synoptic diagram of the proposed solution showing acquisition, filtering, regulation, and transmission/actuation}
    \label{fig:synoptic}
\end{figure}


\section{General Idea Presentation and Used Tools Utility}

As introduced in the previous chapter, the IO‑Box project is conceived as a modular
platform designed to simplify sensor integration and signal management in industrial
contexts. The general idea is to build a system that acquires inputs, processes them
locally, and communicates results to external supervisory platforms. To achieve this, the
IO‑Box relies on both a layered software architecture and a robust hardware foundation.

\subsubsection*{Software Layers}
The software architecture is organized into four distinct layers:
\begin{itemize}
    \item \textbf{Hardware Abstraction Layer (HAL)}: Provides standardized access to
    peripherals and hardware resources, ensuring portability and independence from
    the underlying board.
    \item \textbf{Board Support Package (BSP)}: Contains drivers and low‑level routines
    specific to the STM32 microcontroller and associated components.
    \item \textbf{Manager Layer}: Supervises system functions, including protocol
    management, error handling, and data integrity.
    \item \textbf{Application Layer}: Implements the functional logic of the IO‑Box,
    such as signal processing, actuator control, and communication with external
    supervisory systems.
\end{itemize}

\begin{figure}[H]
    \centering
    \includegraphics[width=0.65\textwidth]{C:/Users/selim/Downloads/imagesreport/software_layers.png}
    \caption{Software architecture organized into four distinct layers}
    \label{fig:software_layers}
\end{figure}

\subsubsection*{Hardware Components}
Complementing the software layers, the hardware architecture provides the physical
interface with the industrial environment. As shown in the block diagrams provided by
the hardware team, the system integrates:
\begin{itemize}
    \item \textbf{Power Supply Architecture}: A multi‑stage conversion system that
    regulates input voltages (12–72V) down to stable levels (48V, 24V, 12V, 5V, and
    3.3V) using buck converters and LDO regulators, ensuring reliable energy
    distribution.
    \item \textbf{Microcontroller Unit (STM32G474)}: Serves as the core processing
    element, coordinating peripherals and executing firmware.
    \item \textbf{Communication Interfaces}: Support for UART, SPI, I\textsuperscript{2}C,
    CAN FD, LIN, and RS485, enabling interoperability with diverse sensors and
    industrial devices.
    \item \textbf{Peripheral Modules}: ADCs, DACs, PWM outputs, multiplexers, and
    transceivers that extend the IO‑Box’s connectivity and control capabilities.
\end{itemize}

Together, these software layers and hardware components form a coherent system that
translates the project’s objectives into a practical and scalable solution. This general
concept provides the foundation for the detailed technical environment and architecture
that will be analyzed in the following sections.

\section{Hardware Tools}
\subsection{STM32G474}

At the heart of the IO‑Box hardware architecture lies the STM32G474 microcontroller,
a high‑performance device from STMicroelectronics. It is based on the ARM Cortex‑M4
core with floating‑point unit (FPU), making it particularly well‑suited for real‑time signal
processing and control applications.

\begin{figure}[H]
    \centering
    \includegraphics[width=0.65\textwidth]{C:/Users/selim/Downloads/imagesreport/stm32g474_chip.jpg}
    \caption{STM32G474 Microcontroller (physical package)}
    \label{fig:stm32_chip}
\end{figure}

The STM32G474 was selected for several reasons:
\begin{itemize}
    \item \textbf{Processing Capabilities}: Equipped with a 170 MHz CPU and DSP
    instructions, it enables efficient execution of filtering, regulation, and control
    algorithms.
    \item \textbf{Rich Peripheral Set}: Provides multiple ADC channels, DAC outputs,
    PWM timers, and communication interfaces (UART, SPI, I\textsuperscript{2}C, CAN FD,
    LIN, USB), ensuring compatibility with diverse sensors and actuators.
    \item \textbf{Memory Resources}: Integrates up to 128 KB SRAM and 512 KB Flash,
    allowing the implementation of complex firmware while maintaining flexibility.
    \item \textbf{Analog Performance}: Features high‑resolution ADCs and DACs,
    supporting precise acquisition and generation of industrial signals.
    \item \textbf{Scalability and Ecosystem}: Benefits from the STM32Cube ecosystem,
    including HAL libraries, BSP drivers, and development tools that simplify firmware
    design and portability.
\end{itemize}

\begin{table}[H]
    \centering
    \caption{Product Attributes — STM32G474 Microcontroller}
    \label{tab:stm32g474_attributes}
    \begin{tabular}{|l|l|}
        \hline
        \textbf{Category} & Integrated Circuits, Embedded Systems, Microcontrollers \\
        \hline
        \textbf{Manufacturer} & STMicroelectronics \\
        \hline
        \textbf{Series} & STM32G4 \\
        \hline
        \textbf{Packaging} & Tray \\
        \hline
        \textbf{Component Status} & Active \\
        \hline
        \textbf{Processor Core} & ARM Cortex‑M4F \\
        \hline
        \textbf{Core Size} & 32‑Bit \\
        \hline
        \textbf{Speed} & 170 MHz \\
        \hline
        \textbf{Connectivity} & CANbus, I\textsuperscript{2}C, IrDA, LINbus, QSPI, SPI, UART/USART \\
        \hline
        \textbf{Peripherals} & DMA, I2S, POR, PWM, WDT, Brown‑out detection/reset \\
        \hline
        \textbf{Number of I/O} & 52 \\
        \hline
        \textbf{Program Memory Size} & 512 KB Flash \\
        \hline
        \textbf{RAM Size} & 128 KB \\
        \hline
        \textbf{Data Converters} & ADC: 26 × 12‑bit; DAC: 7 × 12‑bit \\
        \hline
        \textbf{Supply Voltage} & 1.71 V – 3.6 V \\
        \hline
        \textbf{Oscillator Type} & Internal \\
        \hline
        \textbf{Operating Temperature} & –40 °C to +125 °C \\
        \hline
        \textbf{Mounting Type} & Surface Mount \\
        \hline
        \textbf{Package} & 64‑LQFP (10 × 10 mm) \\
        \hline
        \textbf{Base Product Number} & STM32G474 \\
        \hline
    \end{tabular}
\end{table}



Within the IO‑Box, the STM32G474 acts as the central processing unit, coordinating
signal acquisition, local data processing, and actuator control. It also supervises
communication with external systems through multiple industrial protocols, making it a
key enabler of the project’s objectives.

\section{Software Tools}

\subsection{Visual Studio Code (VS Code)}

Visual Studio Code is employed as a lightweight and versatile integrated development
environment (IDE). It provides syntax highlighting, debugging extensions, and Git
integration, making it suitable for rapid code editing and version control. Its modularity
allows developers to adapt the environment to specific project needs.

\begin{figure}[H]
    \centering
    
\includegraphics[width=0.65\textwidth]{C:/Users/selim/Downloads/imagesreport/vs_code.jpeg}
    \caption{Visual Studio Code interface}
    \label{fig:vscode}
\end{figure}

\subsection{IAR Embedded Workbench}

IAR Embedded Workbench is used for compiling and debugging embedded C code on
the STM32 platform. It offers advanced optimization features, precise debugging
capabilities, and support for hardware breakpoints, which are essential for real‑time
embedded applications. This tool ensures reliable firmware execution on the target
microcontroller.

\begin{figure}[H]
    \centering
    
\includegraphics[width=0.45\textwidth]{C:/Users/selim/Downloads/imagesreport/iar_workbench.png}
    \caption{IAR Embedded Workbench IDE}
    \label{fig:iar}
\end{figure}

\subsection{STM32CubeMX}

STM32CubeMX is a configuration tool provided by STMicroelectronics. It simplifies the
initialization of peripherals, clock trees, and middleware by generating initialization code
automatically. This reduces development time and ensures consistency in hardware
resource management. It also integrates seamlessly with HAL and BSP layers, aligning
with the modular architecture of the IO‑Box.

\begin{figure}[H]
    \centering
    \includegraphics[width=0.65\textwidth]{C:/Users/selim/Downloads/imagesreport/stm32cubemx.png}
    \caption{STM32CubeMX configuration interface}
    \label{fig:cubemx}
\end{figure}

\subsection{GitLab}

GitLab is employed as a collaborative platform for version control and project
management. It enables developers to share code, track changes, and maintain
repositories securely. Through GitLab, the team ensures transparency, reproducibility,
and efficient collaboration, which are critical for academic and industrial projects alike.

\begin{figure}[H]
    \centering
    
\includegraphics[width=0.35\textwidth]{C:/Users/selim/Downloads/imagesreport/gitlab_logo.png}
    \caption{GitLab logo}
    \label{fig:gitlab}
\end{figure}

\section{Communication Protocols}

\subsection{Serial Peripheral Interface (SPI)}

The Serial Peripheral Interface (SPI) is a synchronous serial communication protocol
widely used in embedded systems. It enables high‑speed data exchange between a
microcontroller and peripheral devices such as sensors, memory modules, or
communication transceivers. SPI operates in a master–slave configuration, where the
microcontroller typically acts as the master, controlling the clock and data flow.

In the IO‑Box project, SPI is employed to ensure reliable and efficient communication
with external modules that require fast data transfer. Its simplicity and speed make it
suitable for real‑time applications, particularly when multiple peripherals must be
interfaced with the STM32 microcontroller.

SPI communication is based on a clock signal generated by the master device. Each bit
of data is shifted out on the MOSI line while simultaneously another bit is shifted in on
the MISO line. This process is synchronized by the clock (SCLK), ensuring deterministic
timing. The communication sequence typically follows these steps: the master activates
the Chip Select (CS) line of the desired slave, provides the clock signal, transmits data
bit‑by‑bit on MOSI while receiving data on MISO, samples the data according to the
configured mode, and finally deactivates the CS line to end the session. This mechanism
allows full‑duplex communication, meaning data can be transmitted and received
simultaneously.

SPI supports four modes of operation (Mode 0 to Mode 3), determined by the
configuration of clock polarity (CPOL) and clock phase (CPHA). This flexibility ensures
compatibility with a wide range of peripheral devices. Key characteristics of SPI include
its high speed (operating at several MHz), simple wiring with four main signals (MOSI,
MISO, SCLK, CS), and scalability through multiple slave connections.

\begin{figure}[H]
    \centering
    \includegraphics[width=0.65\textwidth]{C:/Users/selim/Downloads/imagesreport/spi_diagram.png}
    \caption{Basic SPI communication between master and slave devices}
    \label{fig:spi}
\end{figure}

\begin{figure}[H]
    \centering
    \includegraphics[width=0.75\textwidth]{C:/Users/selim/Downloads/imagesreport/endianness_diagram.png}
    \caption{Comparison of Big‑endian and Little‑endian byte ordering in memory}
    \label{fig:endianness}
\end{figure}

\subsection{Controller Area Network (CAN)}

The Controller Area Network (CAN) is a robust serial communication protocol originally
developed for automotive applications, but now widely used in industrial and embedded
systems. It enables reliable data exchange between multiple devices without requiring a
central host, making it ideal for distributed control systems.

In the IO‑Box project, CAN is employed to ensure communication with industrial
equipment and sensors that demand high reliability and fault tolerance. Its ability to
handle multiple nodes on a single bus makes it particularly suitable for environments
where numerous devices must share information in real time.

CAN communication is based on a message‑oriented protocol. Instead of addressing
specific devices directly, each message carries an identifier that defines its priority and
content. This allows multiple nodes to listen simultaneously and decide whether to act
on the received data. The arbitration mechanism ensures that when two nodes attempt
to transmit at the same time, the message with the highest priority (lowest identifier
value) wins, while the other node waits until the bus is free. This guarantees
deterministic communication without collisions.

Key characteristics of CAN include:
\begin{itemize}
    \item \textbf{Multi‑master capability}: Any node can initiate communication when the
    bus is idle.
    \item \textbf{Error detection and handling}: Built‑in mechanisms such as CRC checks,
    acknowledgment bits, and error frames ensure data integrity.
    \item \textbf{Priority‑based arbitration}: Messages are prioritized by identifiers,
    ensuring that critical data is transmitted first.
    \item \textbf{Scalability}: Supports up to 112 nodes on a single bus, making it suitable
    for complex systems.
    \item \textbf{Speed}: Operates up to 1 Mbps in classical CAN, and higher in CAN FD
    (Flexible Data‑Rate), which extends payload size and transmission speed.
\end{itemize}

\begin{figure}[H]
    \centering
    \includegraphics[width=0.70\textwidth]{C:/Users/selim/Downloads/imagesreport/trame_can.jpg}
    \caption{Basic CAN frame}
    \label{fig:can}
\end{figure}


In practice, CAN uses two differential lines, CAN\_H and CAN\_L, to transmit data. This
differential signaling improves noise immunity and allows reliable communication even
in electrically harsh environments. Each frame consists of fields such as the start of
frame, identifier, control bits, data payload, CRC, and end of frame, ensuring structured
and error‑resistant communication.

\begin{figure}[H]
    \centering
    \includegraphics[width=0.70\textwidth]{C:/Users/selim/Downloads/imagesreport/can_diagram.png}
    \caption{Basic CAN bus communication with multiple nodes}
    \label{fig:can}
\end{figure}

\section{Peripherals} 

\subsection{Analog to Digital Converter (ADC)}

The Analog‑to‑Digital Converter (ADC) is a fundamental peripheral in embedded
systems, responsible for transforming continuous analog signals into discrete digital
values that can be processed by the microcontroller. This conversion is essential for
applications where physical phenomena such as temperature, pressure, or current must
be measured and interpreted by digital logic.

In the IO‑Box project, ADCs are employed to acquire signals from sensors and external
modules, enabling precise monitoring and control. The STM32G474 microcontroller
integrates multiple high‑resolution ADC channels, allowing simultaneous sampling of
different inputs. This capability is particularly important in industrial contexts where
multiple signals must be captured in real time.

The operation of an ADC follows a well‑defined process. The analog input is sampled at
regular intervals determined by the sampling frequency. Each sample is then quantized
into a digital value according to the converter’s resolution (e.g., 12‑bit or 16‑bit). The
resulting digital word represents the amplitude of the analog signal at that instant. The
accuracy of this conversion depends on factors such as resolution, sampling rate, and
reference voltage stability.

Key characteristics of ADCs include:
\begin{itemize}
    \item \textbf{Resolution}: Defines the number of discrete levels available.
    \item \textbf{Sampling Rate}: Determines how frequently the analog signal is measured,
    affecting the fidelity of dynamic signals.
    \item \textbf{Reference Voltage}: Sets the maximum measurable input range, ensuring
    proper scaling of the digital output.
    \item \textbf{Input Channels}: Multiple channels allow acquisition from several sensors
    simultaneously.
    \item \textbf{Conversion Modes}: Support for single, continuous, and DMA‑driven
    conversions, enabling flexible integration into firmware.
\end{itemize}

Within the IO‑Box, ADCs are used to capture sensor data such as voltage levels, current
measurements, or other analog signals. These values are then processed by the
firmware to generate meaningful information, trigger control actions, or communicate
results to external supervisory systems.

\begin{figure}[H]
    \centering
    \includegraphics[width=0.70\textwidth]{C:/Users/selim/Downloads/imagesreport/ADC.png}
    \caption{Basic ADC conversion from analog input to digital output}
    \label{fig:adc}
\end{figure}

\subsection{Timers}

Timers are fundamental peripherals in the STM32G474 microcontroller, designed to
provide accurate time base generation, event counting, and signal measurement. They
operate by incrementing or decrementing a counter register according to a clock source,
and can trigger actions when predefined values are reached.

Key functionalities of timers include:
\begin{itemize}
    \item \textbf{Time Base Generation}: Creating precise delays or periodic events.
    \item \textbf{Input Capture}: Measuring external signal characteristics such as
    frequency or pulse width.
    \item \textbf{Output Compare}: Generating events when the counter matches a
    reference value.
    \item \textbf{Synchronization}: Coordinating multiple timers for complex control
    tasks.
    \item \textbf{Advanced Features}: Dead‑time insertion, break inputs, and
    complementary outputs for safe power electronics control.
\end{itemize}

Timers are essential for scheduling tasks, measuring sensor signals, and driving
peripherals that require precise timing.


\subsection{Pulse Width Modulation (PWM)}

Pulse Width Modulation (PWM) is a technique used to encode analog information into a
digital waveform by varying the duty cycle of a periodic signal. In embedded systems,
PWM is widely used for motor control, LED dimming, and power regulation.

PWM signals in the STM32G474 are generated using timers. The timer counts up to a
predefined value stored in the auto‑reload register (ARR), while the compare register
(CCR) defines the switching point of the output. The ratio between the high time and the
total period represents the duty cycle. By adjusting ARR and CCR, both frequency and
duty cycle can be controlled with high precision.

Key characteristics of PWM include:
\begin{itemize}
    \item \textbf{Frequency}: Determined by prescaler and ARR configuration.
    \item \textbf{Duty Cycle}: Defined by CCR relative to ARR.
    \item \textbf{Resolution}: Limited by timer bit width (typically 16‑bit).
    \item \textbf{Complementary Outputs}: Enable safe control of MOSFET drivers in
    half‑bridge or full‑bridge circuits.
    \item \textbf{Applications}: Motor drivers, LED dimming, switching converters,
    and DAC emulation.
\end{itemize}

\begin{figure}[H]
    \centering
    \includegraphics[width=0.70\textwidth]{C:/Users/selim/Downloads/imagesreport/Duty_Cycle_Examples.png}
    \caption{Illustration of PWM signals with different duty cycles (25\%, 50\%, and 75\%)}
    \label{fig:pwm}
\end{figure}

\subsection{EEPROM Memory: M95M01-RDW6TP}

To complement the STM32 microcontroller and ensure reliable data storage, the IO‑Box
integrates an external EEPROM memory, specifically the M95M01‑RMN6TP from STMicroelectronics.
This device provides non‑volatile storage for configuration parameters, calibration data,
and system logs, allowing the platform to retain critical information even after power
cycles.

The M95M01‑RMN6TP is a 1‑Mbit serial EEPROM organized as 131,072 × 8 bits. It
features a fast access time of 25 ns and communicates via the SPI bus, ensuring efficient
integration with the STM32G474. Packaged in an 8‑pin SOIC format, it offers a compact
solution suitable for embedded applications. Its endurance of up to one million write
cycles and data retention of 40 years make it a robust choice for industrial environments
where reliability is essential.

\begin{table}[H]
    \centering
    \caption{Key Characteristics of M95M01‑RDW6TP EEPROM}
    \label{tab:eeprom}
    \begin{tabular}{|l|l|}
        \hline
        \textbf{Feature} & \textbf{Specification} \\
        \hline
        Manufacturer & STMicroelectronics \\
        \hline
        Type & Serial EEPROM \\
        \hline
        Capacity & 1 Mbit (131,072 × 8 bits) \\
        \hline
        Interface & SPI \\
        \hline
        Access Time & 25 ns \\
        \hline
        Package & SOIC‑8 \\
        \hline
        Endurance & 1,000,000 write cycles \\
        \hline
        Data Retention & 40 years \\
        \hline
        Operating Voltage & 1.8 V – 5.5 V \\
        \hline
    \end{tabular}
\end{table}

\begin{figure}[H]
    \centering
    \includegraphics[width=0.55\textwidth]{C:/Users/selim/Downloads/imagesreport/eeprom.png}
    \caption{M95M01‑RDW6TP EEPROM memory chip}
    \label{fig:eeprom}
\end{figure}


\section{Sensors}

In the IO‑Box project, several sensors were integrated to enable environmental and
physical measurements. Our work focused not only on interfacing these sensors with
the STM32G474 microcontroller, but also on configuring their registers, managing
communication protocols, and performing calibration to ensure accurate data
acquisition. Each sensor was studied in detail, and its characteristics were documented
to highlight its role within the system.

\subsection{ENS210 Sensor}

The ENS210 is a digital sensor from ScioSense that integrates both temperature and
humidity measurement in a compact package. It communicates via an I\textsuperscript{2}C
interface, making it easy to integrate with microcontrollers such as the STM32G474.
Its high accuracy and low power consumption make it suitable for environmental
monitoring in embedded systems.

\begin{table}[H]
    \centering
    \caption{Key Characteristics of ENS210 Sensor}
    \label{tab:ens210}
    \begin{tabular}{|l|l|}
        \hline
        \textbf{Feature} & \textbf{Specification} \\
        \hline
        Manufacturer & ScioSense \\
        \hline
        Type & Temperature + Humidity Sensor \\
        \hline
        Interface & I\textsuperscript{2}C (16‑bit output registers) \\
        \hline
        Supply Voltage & 1.71 V – 3.6 V \\
        \hline
        Temperature Range & –40 °C to +100 °C \\
        \hline
        Temperature Accuracy & ±0.15 °C \\
        \hline
        Humidity Range & 0 – 100 \% RH \\
        \hline
       Humidity Accuracy & $\pm$2 \% RH \\
        \hline
        Package & 4‑QFN (2 × 2 mm) \\
        \hline
    \end{tabular}
\end{table}

\begin{figure}[H]
    \centering
    \includegraphics[width=0.55\textwidth]{C:/Users/selim/Downloads/imagesreport/ES-12233_001__04637.jpg}
    \caption{ENS210 digital temperature and humidity sensor}
    \label{fig:ens210}
\end{figure}

\subsection{OPT3001DNPR Sensor}

The OPT3001DNPR is a digital ambient light sensor from Texas Instruments. It is
designed to measure the intensity of visible light and provide results in lux units,
making it ideal for applications such as automatic brightness control, display backlight
adjustment, and environmental monitoring. The sensor communicates via the
I\textsuperscript{2}C protocol and features built‑in calibration and linear response across
a wide dynamic range.

In the IO‑Box project, the OPT3001DNPR was integrated to provide accurate light
measurements. Work was carried out on its internal registers to configure measurement
modes, thresholds, and calibration routines, ensuring reliable data acquisition for
real‑time applications.

\begin{table}[H]
    \centering
    \caption{Key Characteristics of OPT3001DNPR Sensor}
    \label{tab:opt3001}
    \begin{tabular}{|l|l|}
        \hline
        \textbf{Feature} & \textbf{Specification} \\
        \hline
        Manufacturer & Texas Instruments \\
        \hline
        Type & Ambient Light Sensor \\
        \hline
        Interface & I\textsuperscript{2}C \\
        \hline
        Supply Voltage & 1.6 V – 3.6 V \\
        \hline
        Measurement Range & 0.01 lux – 83,000 lux \\
        \hline
        Output Format & 16‑bit digital result in lux \\
        \hline
        Accuracy & ±10\% typical \\
        \hline
        Operating Temperature & –40 °C to +85 °C \\
        \hline
        Package & 6‑Pin WSON (2 × 2 mm) \\
        \hline
    \end{tabular}
\end{table}

\begin{figure}[H]
    \centering
    \includegraphics[width=0.55\textwidth]{C:/Users/selim/Downloads/imagesreport/opt3001.jpg}
    \caption{OPT3001DNPR ambient light sensor}
    \label{fig:opt3001}
\end{figure}






\subsection{Conclusion}

In this chapter, we presented the general concept and system components of the IO‑Box
project. We began with the synoptic diagram, which outlined the main functional phases
of acquisition, filtering, regulation, and transmission/actuation. We then introduced the
software architecture, organized into four distinct layers (HAL, BSP, Manager, and
Application), and complemented this with the hardware architecture, highlighting the
STM32G474 microcontroller, power supply design, communication interfaces, and
peripheral modules. Finally, we described the development tools and protocols that
support the implementation of the system.

By establishing this technical foundation, Chapter 3 has clarified the resources and
methodologies that will be employed in the realization phase. The next chapter will be
dedicated to the practical implementation of the IO‑Box, detailing the firmware design,
register configuration, BSP development, and integration of sensors and actuators into
a coherent and functional platform.

